% UTF8 表示使用通用国际字符作为内码，ctexart 表示这是一篇中文论文
\documentclass[UTF8]{ctexart} 
% 通过调用宏包 geometry 调整页面设置为 A4 纸，页边距 2.6 cm
\usepackage[a4paper,margin=2.6cm]{geometry}

\usepackage{amsmath, amssymb, amsthm} % 数学公式、符号、定理
\usepackage{graphicx}                 % 插图
\usepackage[colorlinks=true,linkcolor=blue]{hyperref} % 超链接与交叉引用

% 设置文章标题
\title{复变函数学习笔记：3.1节“零点与极点”}
% 设置撰写日期
\date{}

% 设置环境 theorem (这是用户自定义环境名，相当于一个变量)
\newtheorem{theorem}{定理}[section]
% 设置引理环境，直接套用theorem，后果是和定理共享编号
\newtheorem{lemma}[theorem]{引理}
% 设置定义环境，和定理共享编号
\newtheorem{definition}[theorem]{定义}
% 设置注记环境，独立编号，和节绑定
\newtheorem{remark}{注记}[section]
% 让公式编号的计数器前面自动加上节，比如式 (2.1)
\numberwithin{equation}{section}

% 自定义命令，专门用于实数域和复数域
\newcommand{\RR}{\mathbb{R}}
\newcommand{\CC}{\mathbb{C}}

\begin{document}

\maketitle

\section{逻辑脉络 (Logical Flow)}
本节的核心目的是为下一节的“留数定理”做局部（Local）分析的准备。其逻辑脉络非常自然且层层递进：
\begin{enumerate}
    \item \textbf{引入孤立奇点}：首先定义孤立奇点，并通过直观例子（如 $z, 1/z, e^{1/z}$）展示奇点的三种不同行为（可去奇点、极点、本性奇点） 。本节主要聚焦于前两类。
    \item \textbf{刻画零点的局部结构}：因为“极点”本质上是“零点”的倒数，教材先分析了全纯函数的零点。通过幂级数展开证明了非平凡全纯函数的零点是\textbf{孤立的}，并提取出“零点阶数（Multiplicity）”的概念 [cite: 158]。
    \item \textbf{刻画极点的局部结构}：利用 $1/f$ 将极点转化为零点的问题，进而得出极点附近函数的标准局部表达形式 $f(z) = (z-z_0)^{-n}h(z)$ [cite: 159]。
    \item \textbf{引入主部分与留数}：将极点的表达式展开，剥离出包含负幂次的“主部分（Principal part）” [cite: 160]。在主部分中，只有 $(z-z_0)^{-1}$ 这一项在闭曲线积分中无法消去，因此它的系数 $a_{-1}$ 被定义为核心概念——\textbf{留数（Residue）} [cite: 160]。
    \item \textbf{推导留数计算公式}：最后，利用求导工具给出了计算 $n$ 阶极点留数的代数公式 [cite: 161]，为后续利用留数定理计算复积分铺平了道路。
\end{enumerate}

\section{主要内容梳理 (Main Content)}

\subsection{孤立奇点与零点}
如果函数 $f$ 在 $z_0$ 的去心邻域内全纯，但在 $z_0$ 处未定义，则称 $z_0$ 为 $f$ 的\textbf{点奇点 (Point singularity)} 或 \textbf{孤立奇点 (Isolated singularity)} 。

对于全纯函数的零点，教材给出了如下极其重要的局部结构定理：

\begin{theorem}
假设 $f$ 在连通开集 $\Omega$ 上全纯，
在 $z_0 \in \Omega$ 处有一个零点，
且 $f$ 在 $\Omega$ 上不恒为零。
那么存在 $z_0$ 的一个邻域 $U \subset \Omega$，
一个在 $U$ 上\textbf{无零点}的全纯函数 $g$，以及一个唯一的正整数 $n$，使得对所有 $z \in U$ 有：
$$ f(z) = (z-z_0)^n g(z) $$ [cite: 158]
\end{theorem}

\begin{remark}
这里的整数 $n$ 被称为零点的\textbf{阶 (order / multiplicity)} [cite: 159]。该定理证明了非平凡全纯函数的\textbf{零点必定是孤立的} [cite: 158]。
\end{remark}

\subsection{极点与主部分}
对于孤立奇点 $z_0$，如果定义 $1/f(z_0) = 0$ 后，$1/f$ 可以在 $z_0$ 的一个完整邻域内全纯，则称 $z_0$ 为 $f$ 的\textbf{极点 (Pole)} [cite: 159]。基于零点的定理，可以直接推导出极点的结构：

\begin{theorem}
如果 $f$ 在 $z_0$ 处有极点，那么在该点的邻域内，存在一个无零点的全纯函数 $h$ 和一个唯一的正整数 $n$，使得：
$$ f(z) = (z-z_0)^{-n} h(z) $$ [cite: 159]
\end{theorem}
这里的 $n$ 称为极点的阶，它描述了函数在 $z_0$ 附近趋于无穷的速度。如果 $n=1$，称为简单极点 (simple pole) [cite: 160]。

将 $h(z)$ 在 $z_0$ 处展为泰勒级数并代入上式，可以得到极点附近的展开式：
\begin{theorem}
如果 $f$ 在 $z_0$ 处有 $n$ 阶极点，则：
$$ f(z) = \frac{a_{-n}}{(z-z_0)^n} + \frac{a_{-n+1}}{(z-z_0)^{n-1}} + \dots + \frac{a_{-1}}{z-z_0} + G(z) $$
其中 $G(z)$ 在 $z_0$ 的邻域内全纯 [cite: 160]。
\end{theorem}

\begin{proof}
    $ f(z) = (z-z_0)^{-n} h(z) $,h是全纯的,
    故对$h(z)$进行幂级数展开,并且带入这个式子即可证明
\end{proof}

\begin{definition}
上式中的负幂次部分 $\sum_{k=1}^n a_{-k}(z-z_0)^{-k}$ 称为 $f$ 在 $z_0$ 处的\textbf{主部分 (Principal part)} [cite: 160]。系数 $a_{-1}$ 称为 $f$ 在该极点处的\textbf{留数 (Residue)}，记作 $\mathrm{res}_{z_0} f = a_{-1}$ [cite: 160]。
\end{definition}

\subsection{留数的计算公式}
留数之所以特殊，是因为在主部分中，除了 $(z-z_0)^{-1}$ 之外，其余高阶负幂次项在去心邻域内都存在原函数 [cite: 160]。因此，如果 $C$ 是围绕 $z_0$ 的圆周，则只有 $a_{-1}$ 会对积分产生非零贡献：$\frac{1}{2\pi i}\int_C P(z)dz = a_{-1}$ [cite: 160]。

为了便于实际计算，教材给出了高阶极点留数的计算公式：
\begin{theorem}
如果 $f$ 在 $z_0$ 处有 $n$ 阶极点，则其留数可以通过下式计算：
$$ \mathrm{res}_{z_0} f = \lim_{z \to z_0} \frac{1}{(n-1)!} \left(\frac{d}{dz}\right)^{n-1} \left[ (z-z_0)^n f(z) \right] $$ [cite: 161]
\end{theorem}
特别地，当 $z_0$ 是简单极点 ($n=1$) 时，公式退化为：$\mathrm{res}_{z_0} f = \lim_{z \to z_0} (z-z_0)f(z)$ [cite: 160]。

\end{document}